\chapter{从已有数据自己回测一天}

\section{本章要解决的困惑}

本章把前面内容变成一次新人可以执行的单日路线。重点不是背命令，而是知道每条命令验证哪条链路。

\section{第零步：先选一个非常小的目标}

新人不要一上来跑三周、多 profile、多币种。第一天只做：

\begin{lstlisting}
one symbol
one date
one policy
one admission profile
one exit profile
\end{lstlisting}

目标也不是找新 alpha，而是确认自己能从输入数据解释输出 artifact。一个合格的单日实验，哪怕结果是亏的，也应该能回答：

\begin{quote}
这一天哪些 frame 触发了？哪些被 spread gate 拒绝？哪些被 3x ledger 剪掉？60 秒 executable 是怎样算的？
\end{quote}

\section{第一步：检查 decision frame cache}

\begin{lstlisting}
python systems/ccusdt_replay_exchange/diagnostics/decision_frame_cache.py validate ^
  --repo-root . ^
  --symbol CCUSDT ^
  --from-date 2026-05-18 ^
  --to-date 2026-05-18
\end{lstlisting}

它验证：

\begin{quote}
canonical trade/L2 是否已经被转换成 runtime-safe 的 \code{decision_frame_v1}。
\end{quote}

\section{第二步：理解你要跑的 profile}

跑之前先写清楚：

\begin{lstlisting}
symbol = CCUSDT
date = 2026-05-18
policy = old_tfi_four_cell_shadow_v1
admission = entry_spread_q70_v1
capacity = core_idle01 or fifo_clip
exit = fixed60_taker or fixed60_mid
leverage_cap = 3
\end{lstlisting}

如果 profile 写不清，输出结果也没有解释力。

\section{第三步：跑 fast backtest}

示意命令：

\begin{lstlisting}
python systems/ccusdt_replay_exchange/diagnostics/fast_strategy_backtest.py ^
  --repo-root . ^
  --symbol CCUSDT ^
  --from-date 2026-05-18 ^
  --to-date 2026-05-18 ^
  --exit-profile fixed60_taker ^
  --capacity-profile core_idle01 ^
  --admission-profile entry_spread_q70_v1 ^
  --leverage-cap 3
\end{lstlisting}

实际项目中 q70 admission 可能还需要 threshold JSON。没有阈值不要假装等价。

\section{第四步：检查输出}

顺序：

\begin{lstlisting}
summary.json
profile_manifest.json
daily.csv
entries.parquet
exits.parquet
\end{lstlisting}

先问：

\begin{itemize}[leftmargin=2em]
  \item actual entries 多少？
  \item admission 拒绝了多少？
  \item capacity clipped 多少？
  \item mid 和 executable 差多少？
  \item daily 是否稳定？
\end{itemize}

\section{读一行 entries.parquet}

假设你从 \code{entries.parquet} 抽出一行：

\begin{lstlisting}
entry_ts_us = 1716000010000000
side = buy
cell = 11_r5_frames
membership_set = tfi_follow_flat+tfi_long_flat+tfi_event_active
entry_bid = 0.1120
entry_ask = 0.1124
entry_spread_bps = 35.65
entry_cross_bps = 17.825
a_r5_raw = 4.33
base_weight = 0.5
gamma = 0.75
requested_exposure = 0.375
actual_exposure = 0.1
\end{lstlisting}

你应该能逐项解释：

\begin{enumerate}[leftmargin=2em]
  \item \code{side=buy} 来自 \(TFI=1\)；
  \item \code{entry_cross_bps} 是 half spread；
  \item \code{cell=11_r5_frames} 来自 \(R5\ge1\) 且 \(frames\ge31\)；
  \item \code{requested_exposure=0.375} 来自 base weight 乘 gamma；
  \item \code{actual_exposure=0.1} 表示 3x ledger 只剩 0.1 free capacity。
\end{enumerate}

如果这五项解释不了，说明还没真正掌握当前策略。

\section{读一行 exits.parquet}

假设对应 exit：

\begin{lstlisting}
entry_side = buy
entry_ask = 0.1124
exit_bid = 0.1126
exit_profile = fixed60_taker
net_bps = 17.8
\end{lstlisting}

手算：

\[
net=10000\log\frac{0.1126}{0.1124}=17.8.
\]

再乘 exposure：

\[
contribution=17.8\times0.1=1.78.
\]

这才是它对 weighted result 的贡献。只看 \code{net_bps} 不看 \code{actual_exposure}，会误判策略贡献。

\section{读 daily.csv}

单日 \code{daily.csv} 应该和 entries/exits 能对上：

\begin{longtable}{p{0.28\textwidth}p{0.56\textwidth}}
\toprule
字段 & 手动解释 \\
\midrule
\code{actual_entries} & \code{actual_exposure > 0} 的 entry 数 \\
\code{gross_entries} & policy/admission/capacity 前后的候选口径要看定义 \\
\code{net_weighted_bps} & 每笔 executable bps 乘 actual exposure 后求和 \\
\code{admission_rejected} & spread gate 拒绝数量 \\
\code{capacity_clipped} & requested 大于 free capacity 的数量或总量 \\
\bottomrule
\end{longtable}

不要把 \code{actual_entries} 当成交易量，也不要把 raw trigger count 当成资金投入。仓位贡献要看 \code{actual_exposure}。

\section{第五步：必要时 strict replay}

如果只是研究 entry/gamma，fast 够了。如果要验收 Bot/Runner 路径，用：

\begin{lstlisting}
cargo run --manifest-path systems/ccusdt_replay_exchange/engine/Cargo.toml ^
  -p ccusdt_replay_cli -- run sparse-python ^
  --canonical-date 2026-05-18 ^
  --max-events 500 ^
  --latency-us 50000
\end{lstlisting}

它验证的是：

\begin{quote}
Runner 能否把 market stream 送给 Python Bot，并把 submit order 转成 arrival、fill、portfolio 和 events.ndjson。
\end{quote}

\section{单日实验成功的标准}

不是赚了就成功。成功是你能解释：

\begin{enumerate}[leftmargin=2em]
  \item 每笔 entry 为什么触发；
  \item 每笔 entry 为什么被接受或拒绝；
  \item 每笔 actual exposure 为什么被剪裁；
  \item fixed60 executable 是怎样算的；
  \item fast 结果和 strict 审计分别说明什么。
\end{enumerate}

\checkpoint{自己回测的第一目标不是赚钱，而是能解释每一行输出。}
